% Part: proof-theory
% Chapter: proof-search
% Section: completeness

\documentclass[../../../include/open-logic-section]{subfiles}

\begin{document}

\olfileid{pt}{ps}{cpl}

\olsection{بشپړتيا}

اوس ښيو چې د ثبوت د لټون الګوريتم بې‌خطا دے: يعنې که
درېږي او ناکامېږي يا هېڅکله نۀ درېږي، د پيل
سيکوېنټ~$\Gamma \Sequent \Delta$ هېڅ !!{proof} نشته۔ د دې
دپاره !!a{structure}~$\Struct M$ داسې تعريفوو چې د هر
$!A \in \Gamma$ دپاره $\Sat{M}{!A}$، او د هر~$!A \in \Delta$
دپاره $\Sat/{M}{!A}$ وي۔ په بله وينا، په~$\Struct M$ کښې
د~$\Gamma$ ټول !!{formula}s رښتيا او د~$\Delta$ ټول
!!{formula}s دروغ دي؛ يعنې $\Gamma \Sequent \Delta$ معتبر
نۀ دے۔ ځکه چې \Log{G3c} سم دے، $\Gamma \Sequent \Delta$
هېڅ !!{proof} نۀ لري۔
\textbf{د سرچينې د نښې سمون OLCMP-185:} په ښي اړخ کښې
د فارمول نيمګړې نښه په هماغه بشپړه نښه سمه شوې ده۔

$\Struct{M}$ د !!{formula}s له يوې جوړې $\Theta, \Xi$
او د !!{constant}s له سېټ~$C$ څخه تعريفوو۔ $\Theta, \Xi$ او~$C$
د ناکام شوي، يعنې درېدلي يا نۀ درېدېدونکي، ثبوت لټون له
\emph{ناکامې څانګې} څخه اخلو۔ ناکامه څانګه د سيکوېنټونو
$\Pi_n \Sequent \Lambda_n$ او د !!{constant}s د سېټونو~$C_n$
متناهي يا نامتناهي لړۍ ده، چې $\Pi_0 \Sequent \Lambda_0$
ئې پاييز سيکوېنټ $\Gamma \Sequent \Delta$ دے۔ هر چې راتلونکے
مورد موجود وي، $\Pi_{n+1} \Sequent \Lambda_{n+1}$ د
$\Pi_n \Sequent \Lambda_n$ هغه مقدمه ده چې په پړاو~$n+1$
کښې جوړه شوې او الګوريتم ئې !!a{proof} نۀ مومي۔ د هر~$n$ په داسې
ناتمام پړاو کښې بايد لږ تر لږه يوه داسې مقدمه وي؛ کنه
الګوريتم به !!a{proof} موندلے و۔ $C_n$ د !!{constant}s هغه
سېټ دے چې له~$\Pi_n \Sequent \Lambda_n$ سره تړلے دے۔
\textbf{د سرچينه‌يي څانګې سمون OLCMP-186:} د درېدلې
متناهي څانګې تر وروستي مورد وروسته راتلونکې مقدمه نۀ
غوښتل کېږي؛ وروستۍ کاپي شوې ښي‌خوا نښه هم د همدې څانګې
په اصلي نښه سمه شوې ده۔

که الګوريتم پر هغه تر ټولو پاسني سيکوېنټ ناکام ودرېږي
چې يوازې اټومي !!{formula}s لري، نو په همدغه سيکوېنټ ختمه
څانګه ناکامه څانګه ده۔ که الګوريتم هېڅکله نۀ درېږي، پرله‌پسې
لويې ونې جوړوي۔ دا د يوې نامتناهي ونې وده بللے شئ: هغه ونه
چې د الګوريتم له نامتناهي ډېرو ګامونو وروسته به موجوده وه۔
دا ونه نامتناهي، خو متناهي څانګېدونکې وي؛ هر غوټه يوازې يو
يا دوه اولادونه لري، يعنې هغه سيکوېنټونه چې سمدستي ئې د
پاسه دي۔ د کونيګ د لمې له مخې هره نامتناهي، متناهي
څانګېدونکې ونه يوه نامتناهي څانګه لري۔ که د لټون الګوريتم
تل روان وي، همدغه نامتناهي څانګه ناکامه څانګه ده۔

که د $\Pi_n \Sequent \Lambda_n$ او $C_n$ لړۍ ناکامه څانګه
وي، $\Theta = \bigcup_n \Pi_n$ او $\Xi = \bigcup_n \Lambda_n$
واخلئ، چې شاخصونه او د !!{formula}s تکراري موردونه ترې لرې
کوو؛ همدارنګه $C = \bigcup_n C_n$ واخلئ۔

اوس~$\Struct{M}$ تعريفوو۔
\begin{enumerate}
\item د تعريف ساحه $\Domain{M}$ د هغو ترمونو سېټ دے چې له~$C$
او په~$\Gamma \Sequent \Delta$ کښې له موجودو !!{function}s
څخه جوړېږي، خو هېڅ !!{variable}s نۀ لري۔
\item د هر ثابت $c \in \Domain{M}$ دپاره، $\Assign{c}{M} = c$۔
\item که $f$ په $\Gamma \Sequent \Delta$ کښې يو $m$-ځايه
!!{function} وي، او $t_1$، \dots، $t_m \in \Domain{M}$ وي، نو
$\Assign{f}{M}(t_1, \dots, t_m) = f(t_1, \dots, t_m)$۔
\item که $R$ په $\Gamma \Sequent \Delta$ کښې يو $m$-ځايه
!!{predicate} وي، نو $\Assign{R}{M} = \Setabs{\tuple{t_1, \dots, t_m} \in
\Domain{M}^m}{R(t_1, \dots, t_m) \in \Theta}$۔
\end{enumerate}
$\Struct{M}$ يو \emph{ترمي مدل} دے: د تعريف ساحه ئې د ترمونو
سېټ دے او !!{constant}s او !!{function}s داسې تفسير شوي چې
$\Value{t}{M} = t$ تضمين کړي۔ په~$\Theta$ کښې اټومي
!!{formula}s د $\Assign{R}{M}$ د تعريف دپاره داسې کارېږي چې
$\Sat{M}{R(t_1, \dots, t_n)}$ هله او يوازې هله چې
$R(t_1, \dots, t_n) \in \Theta$ وي۔
\textbf{د سرچينه‌يي تفسير سمون OLCMP-187:} د ثابت تفسير
د ثابت دپاره دے، نۀ د هر مرکب ترم د نوم دپاره۔
\textbf{د ځايونو د شمېر سمون OLCMP-188:} د تعريف په دواړو
ځايونو کښې د اعلان شوي شمېر پر ځاے کاپي شوے بل شاخص په
هماغه اعلان شوي شمېر بدل شوے دے۔

\begin{lem}\ollabel{lem:truth}
د هر $!A \in \Theta \cup \Xi$ دپاره:
\begin{enumerate}
  \item که $!A \in \Theta$ وي، نو $\Sat{M}{!A}$۔
  \item که $!A \in \Xi$ وي، نو $\Sat/{M}{!A}$۔
\end{enumerate}
\textbf{د سرچينې د ثبوت د حد څرګندونه OLCMP-190:} لاندې
بې‌متغيره ترمي جوړونه د تړلو جملو حالت ښيي۔ د ازادو
متغيرونو د ګومارنې يا په نوو ثابتونو د هغو د بدلولو جلا
جوړونه په ورکړل شوي دليل کښې نشته؛ عمومي رسمي جمله عين
ساتل شوې، خو دغه اضافي حالت دلته ثابت نۀ ادعا کېږي۔
\end{lem}

\begin{proof}
  د $\depth{!A}$ پر سر په استقرا۔

لومړے فرض کړئ چې $!A$ اټومي دے؛ يعنې يا
$!A \ident \lfalse$ او يا $!A \ident R(t_1, \dots, t_m)$۔

ځکه چې $\Pi_n \Sequent \Lambda_n$ ناکامه څانګه ده، هېڅ
$\Pi_n$ د $\lfalse$ فارمول نۀ شي لرلے؛ کنه
$\Pi_n \Sequent \Lambda_n$ به بديهيت و۔ د~$\Struct{M}$ د
تعريف له مخې $\Sat{M}{R(t_1, \dots, t_m)}$ هله او يوازې هله
چې $R(t_1, \dots, t_m) \in \Theta$ وي۔ نو دواړه خبرې يوځاے
دا راکوي: که $!A \in \Theta$ وي، نو $\Sat{M}{!A}$۔

که $!A$ اټومي او $!A \in \Pi_n$ وي، نو د راتلونکي
مورد د موجودېدو په صورت کښې $!A \in \Pi_{n+1}$؛ او که
$!A \in \Lambda_n$ وي، نو $!A \in \Lambda_{n+1}$۔ دا د لټون
الګوريتم د راتلونکو سيکوېنټونو د جوړولو له طريقې راځي۔ نو
که $!A \in \Theta \cap \Xi$ وي، د يو~$n$ دپاره
$!A \in \Pi_n \cap \Lambda_n$ وي۔ دا مانا لري چې
$\Pi_n \Sequent \Lambda_n$ بديهيت دے، خو ناکامه څانګه هېڅ
بديهيت نۀ لري۔ نو د اټومي~$!A$ دپاره، د جوړونې له مخې
$!A \notin \Theta \cap \Xi$؛ ځکه که $!A \in \Xi$ وي، نو
$!A \notin \Theta$۔ د~$\Struct{M}$ د تعريف له مخې، له دې
راځي چې که $!A \in \Xi$ وي، نو $\Sat/{M}{!A}$۔

د استقرا د ګام دپاره فرض کړئ چې (۱) او (۲) د هغو ټولو
!!{formula}s دپاره صحيح دي چې !!{depth} ئې له~$!A$ څخه لږه
ده۔ د حالتونو په بېلولو د $!A$ دپاره (۱) او (۲) ثابتوو۔

لومړے د ناعټومي فارمول د ټاکلو انصاف په پام کښې ونيسئ۔
کله چې $!A^i$ په $\Pi_n \Sequent \Lambda_n$ کښې موجود وي،
تر څو ټاکل شوے نۀ وي، $!A^i$ په
$\Pi_{n+1} \Sequent \Lambda_{n+1}$ کښې هم پاتې کېږي۔ ټاکنه
په $\Pi_n \Sequent \Lambda_n$ کښې د ټولو شاخصونو پر ځاے د ناعټومي فارمولونو تر ټولو کوچنے
شاخص~$i$ کاروي۔ د منصفانه لټون له مخې بالاخره داسې
$\Pi_n \Sequent \Lambda_n$ ته رسېږو چې $!A^i$ پکښې موجود
وي او $i$ د ناعټومي فارمولونو تر ټولو کوچنے شاخص وي۔ په
پړاو $n+1$ کښې الګوريتم~$!A^i$ راکموي۔ په بله وينا، که
$!A \in \Theta$ وي، د يو $n$ دپاره $\Pi_n = \Pi_n', !A^i$
وي او $i$ د ناعټومي فارمولونو تر ټولو کوچنے شاخص وي۔ او
که $!A \in \Xi$ وي، د يو $n$ دپاره
$\Lambda_n = \Lambda_n', !A^i$ وي او $i$~ئې همدغه تر ټولو
کوچنے ناعټومي شاخص وي۔
\textbf{د سرچينه‌يي ټاکنې سمون OLCMP-189:} اټومي فارمولونه
نۀ ټاکل کېږي۔ د ټولو شاخصونو تر ټولو کوچنے شاخص د
پاتې اټومي موردونو له کبله په هر پړاو کښې سخت نۀ زياتېږي۔
په الګوريتم کښې د OLCMP-192 تعقيب نوي شاخصونه له اوسني
اعظمي شاخصه پورته ټاکلي؛ دلته د مقدمو هماغه نښې ورسره
همغږې شوې دي۔ د ناعټومي موردونو د ټاکنې انصاف ساتل کېږي؛
د ازادو متغيرونو جلا پاتې عمومي تشه په دې سمون نۀ حل کېږي۔

\begin{enumerate}
  \item $!A \in \Theta$ او $!A \ident !B \land !C$۔ نو د يو
  $n$ دپاره $\Pi_n = \Pi_n', (!B \land !C)^i$۔
  $\Pi_{n+1} \Sequent \Lambda_{n+1}$ د~\LeftR{\land} اړونده
  مقدمه ده؛ يعنې $\Pi_{n+1} = \Pi_n', !B^{k+1}, !C^{k+2}$۔ نو
  $!B, !C \in \Theta$۔ د استقرا د فرضيې له مخې
  $\Sat{M}{!B}$ او $\Sat{M}{!C}$۔ د $\Sat{M}{}$ د تعريف له
  مخې، $\Sat{M}{!B \land !C}$ لرو۔

  \item $!A \in \Xi$ او $!A \ident !B \land !C$۔ نو د يو
  $n$ دپاره $\Lambda_n = \Lambda_n', !B \land !C$۔
  $\Pi_{n+1} \Sequent \Lambda_{n+1}$ د~\RightR{\land} له
  دوو مقدمو څخه يوه ده، چې پايله ئې
  $\Pi_n \Sequent \Lambda'_n, !B \land !C$ ده؛ يعنې
  $\Lambda_{n+1} = \Lambda_n', !B^k$ يا
  $\Lambda_{n+1} = \Lambda_n', !C^k$۔ نو يا $!B \in \Xi$ يا
  $!C \in \Xi$۔ د استقرا د فرضيې له مخې، يا $\Sat/{M}{!B}$
  يا $\Sat/{M}{!C}$۔ د $\Sat{M}{}$ د تعريف له مخې
  $\Sat/{M}{!B \land !C}$ لرو۔

  \item هغه حالتونه چې $!A \ident \lnot !B$،
  $!A \ident !B \lor !C$ او $!A \ident !B \lif !C$ وي، ورته
  دي او د مشقونو په توګه پاتې دي۔

  \item $!A \in \Theta$ او $!A \ident \lexists[x][!B(x)]$۔
  نو د يو $n$ دپاره $\Pi_n = \Pi_n', \lexists[x][!B(x)]^i$۔
  $\Pi_{n+1} \Sequent \Lambda_{n+1}$ د~\LeftR{\lexists}
  اړونده مقدمه ده؛ يعنې د يو~$c$ دپاره
  $\Pi_{n+1} = \Pi_n', !B(c)^{k+1}$۔ نو $!B(c) \in \Theta$ او
  $c \in C$۔ د استقرا د فرضيې له مخې $\Sat{M}{!B(c)}$۔ د
  $\Sat{M}{}$ د تعريف له مخې، $\Sat{M}{\lexists[x][!B(x)]}$
  لرو۔

  \item $!A \in \Xi$ او $!A \ident \lexists[x][!B(x)]$۔ نو
  د يو $n$ دپاره $\Lambda_n = \Lambda_n', \lexists[x][!B(x)]^i$۔
  هر ځل چې $\lexists[x][!B(x)]$ راکمېږي،
  $\lexists[x][!B(x)]$ د مقدمې سيکوېنټ په ښي اړخ کښې له نوي
  شاخص سره پاتې کېږي؛ نو که $\lexists[x][!B(x)]$ په ناکامه څانګه کښې موجود وي،
  په ښي اړخ کښې نامتناهي ځله راکمېږي۔ د ناکامې څانګې پر
  سر هر ترم $t \in \Domain{M}$ بالاخره هغه لومړنے ترم شي
  چې يوازې د~$C_n$ ثابتونه لري او تر دې وړاندې په ښي اړخ
  کښې د $\lexists[x][!B(x)]$ په راکمولو کښې نۀ دے کارېدلے۔
  نو د هر $t \in \Domain{M}$ دپاره، د يو~$n$ لپاره
  $!B(t) \in \Lambda_n$ وي او ځکه $!B(t) \in \Xi$۔ د استقرا
  د فرضيې له مخې، د هر $t \in \Domain{M}$ دپاره
  $\Sat/{M}{!B(t)}$۔ د $\Sat{M}{}$ د تعريف له مخې،
  $\Sat/{M}{\lexists[x][!B(x)]}$ لرو۔

  \item هغه حالتونه چې $!A \ident \lforall[x][!B(x)]$ وي،
  ورته دي او د مشقونو په توګه پاتې دي۔
\end{enumerate}
\end{proof}

\begin{cor}\ollabel{cor:complete} د \Log{G3c} د ثبوت د
لټون الګوريتم بشپړ دے: که ناکام ودرېږي يا هېڅکله نۀ
درېږي، د پيل سيکوېنټ~$\Gamma \Sequent \Delta$ معتبر نۀ دے
او ځکه هېڅ !!{proof} نۀ لري۔
\end{cor}

\begin{proof}
  که الګوريتم ناکام ودرېږي يا هېڅکله نۀ درېږي، يوه ناکامه
  څانګه شته چې~$\Struct{M}$ ترې تعريفېږي۔ د
  \olref{lem:truth} له مخې، د هر $!A \in \Gamma$ دپاره
  $\Sat{M}{!A}$ او د هر $!A \in \Delta$ دپاره
  $\Sat/{M}{!A}$۔ په ياد ولرئ چې $\Pi_0 = \Gamma$ او
  $\Lambda_0 = \Delta$، نو $\Gamma \subseteq \Theta$ او
  $\Delta \subseteq \Xi$۔ ځکه $\Gamma \Sequent \Delta$
  معتبر نۀ دے۔ د~\Log{G3c} د سموالي له مخې،
  $\Gamma \Sequent \Delta$ هېڅ !!{proof} نۀ لري۔
\end{proof}

\begin{cor}
\Log{G3c} بشپړ دے: که $\Gamma \Sequent \Delta$ معتبر وي، نو
$\Log{G3c} \Proves \Gamma \Sequent \Delta$۔
\end{cor}

\begin{proof}
  عکس نقيض ثابتوو۔ فرض کړئ چې $\Log{G3c} \Proves/ \Gamma
  \Sequent \Delta$۔ نو د ثبوت د لټون الګوريتم بايد ناکام
  ودرېږي يا هېڅکله نۀ درېږي، ځکه هېڅ !!{proof} نۀ شته چې
  وموندل شي۔ د \olref{cor:complete} له مخې،
  $\Gamma \Sequent \Delta$ معتبر نۀ دے۔
\end{proof}

\end{document}
